% Consumo e impuestos verdes · Parte 2: impuestos correctivos
% Finanzas Públicas, Maestría en Economía Aplicada, Tec de Monterrey
% Esquema aprobado: CLAUDE.md de esta carpeta (2026-10-12). Compilar con: latexmk -pdf consumo_verdes.tex
% Versión para repartir: consumo_verdes_handout.tex (misma fuente, opción handout).
\documentclass[aspectratio=169,11pt]{beamer}

\usepackage[T1]{fontenc}
\usepackage{lmodern}
\usepackage[spanish,es-nodecimaldot]{babel}
\unaccentedoperators % operadores sin acento: max, min (la opción es-noaccentedoperators no existe en esta versión de babel)
\usepackage{amsmath,amssymb}
\usepackage{booktabs}
\usepackage{tikz}
\usetikzlibrary{babel,arrows.meta,calc,positioning,decorations.pathreplacing,patterns}
\usepackage{pgfplots}
\pgfplotsset{compat=1.18}

% ---------------------------------------------------------------------------
% Tema sobrio (igual al de la sesión del 26 de septiembre)
% ---------------------------------------------------------------------------
\definecolor{cTitulo}{RGB}{0,55,100}
\definecolor{cOferta}{RGB}{31,87,153}
\definecolor{cDemanda}{RGB}{200,100,0}
\definecolor{cImpuesto}{RGB}{175,30,30}
\definecolor{cRecauda}{RGB}{40,130,80}
\definecolor{cGris}{RGB}{110,110,110}

\usetheme{default}
\setbeamertemplate{navigation symbols}{}
\setbeamercolor{frametitle}{fg=cTitulo}
\setbeamercolor{title}{fg=cTitulo}
\setbeamercolor{structure}{fg=cTitulo}
\setbeamercolor{block title}{bg=cTitulo!12,fg=cTitulo}
\setbeamercolor{block body}{bg=cTitulo!4}
\setbeamercolor{block title alerted}{bg=cImpuesto!12,fg=cImpuesto}
\setbeamercolor{block body alerted}{bg=cImpuesto!4}
\setbeamerfont{frametitle}{size=\large,series=\bfseries}
\setbeamertemplate{itemize items}[circle]
\setbeamertemplate{blocks}[rounded]
\setbeamertemplate{footline}{%
  \hfill{\small\color{cGris}Consumo e impuestos verdes · Impuestos verdes\quad\insertframenumber/\inserttotalframenumber}\hspace*{1em}\vspace*{0.6em}}

% Fuente al pie de cada diapositiva
\newcommand{\fuente}[1]{\par\vfill{\small\color{cGris}#1\par}}
% Resultado final de una derivación, resaltado
\newcommand{\resultado}[1]{\colorbox{cTitulo!12}{$\displaystyle #1$}}
% \ejes{xmax}{ymax}{etiqueta x}{etiqueta y}
\newcommand{\ejes}[4]{%
  \draw[eje] (0,0) -- (#1,0) node[below] {#3};
  \draw[eje] (0,0) -- (0,#2) node[left] {#4};}

% ---------------------------------------------------------------------------
% Estilo común de gráficas
% ---------------------------------------------------------------------------
\tikzset{
  every node/.append style={font=\small},
  eje/.style={-{Stealth[length=2mm]},thick,black!75},
  flecha/.style={-{Stealth[length=2mm]},thick,cTitulo},
  guia/.style={thin,dashed,black!45},
  punto/.style={circle,fill=black,inner sep=1.2pt},
  caja/.style={draw=cTitulo,rounded corners,fill=cTitulo!5,align=center,inner sep=4pt},
  cajafisco/.style={draw=cRecauda,rounded corners,fill=cRecauda!10,align=center,inner sep=3pt},
  demanda/.style={very thick,cDemanda},
  oferta/.style={very thick,cOferta},
  social/.style={very thick,cImpuesto},
  cuna/.style={very thick,cImpuesto,{Stealth[length=2mm]}-{Stealth[length=2mm]}},
  gana/.style={fill=cRecauda!35},
  pierde/.style={fill=cImpuesto!35},
}

\title{Consumo e impuestos verdes}
\subtitle{Parte 2: impuestos correctivos}
\author{Héctor Juan Villarreal}
\institute{Finanzas Públicas · Maestría en Economía Aplicada\\Escuela de Gobierno y Transformación Pública · Tec de Monterrey}
\date{Octubre de 2026}


\begin{document}

% =========================================================================
% I. TEORÍA
% =========================================================================
\section{Teoría}

% 1 -----------------------------------------------------------------------
\begin{frame}[plain]
  \titlepage
\end{frame}

% 2 -----------------------------------------------------------------------
\begin{frame}{Cuando el triángulo es ganancia}
  \begin{columns}[T]
    \begin{column}{0.42\textwidth}
      \begin{itemize}
        \item La refinería paga el crudo y los salarios: es su \textbf{costo marginal privado}.
        \item No paga el aire sucio de la colonia de al lado: el \textbf{daño marginal}.
        \item Costo social $=$ privado $+$ daño. El mercado produce $x_m$; lo eficiente es $x^*$.
      \end{itemize}
      \pause
      \begin{block}{Impuesto pigouviano}
        $\tau^*$ = daño marginal en el óptimo. La cantidad cae de $x_m$ a $x^*$ y la sociedad \textbf{gana} el triángulo verde.
      \end{block}
    \end{column}
    \begin{column}{0.56\textwidth}
      \centering
      \begin{tikzpicture}[x=0.95cm,y=0.85cm]
        \onslide<2->{\fill[gana] (3.143,2.8) -- (4.286,3.2) -- (4.286,2.0) -- cycle;}
        \ejes{6.4}{5.6}{$x$}{pesos}
        \draw[demanda] (0,5) -- (6,0.8) node[below,pos=0.97] {$D$};
        \draw[oferta] (0,0.5) -- (6,2.6);
        \node[cOferta,below right] at (0.4,0.64) {CMg privado};
        \draw[social] (0,1.7) -- (6,3.8) node[pos=0.85,above=2pt] {CMg social};
        \draw[guia] (4.286,2.0) -- (4.286,0) (3.143,1.6) -- (3.143,0);
        \node[punto] at (4.286,2.0) {}; \node[punto] at (3.143,2.8) {};
        \node[below] at (4.286,0) {$x_m$}; \node[below] at (3.143,0) {$x^*$};
        \onslide<2->{\draw[cuna] (3.143,1.6) -- (3.143,2.8);
        \node[left,cImpuesto] at (3.143,2.2) {$\tau^*$};}
      \end{tikzpicture}
    \end{column}
  \end{columns}
  \fuente{Gruber, caps.~5.1 y 5.3; gráfica propia. Contraste: sesión del 26 de septiembre, diap.~37.}
\end{frame}

% 3 -----------------------------------------------------------------------
\begin{frame}{Derivación mínima}
  Un bien $x$ con beneficio $B(x)$, costo $C(x)$ y daño a terceros $D(x)$.
  \begin{align*}
    \onslide<2->{\text{Bienestar social:}\quad & W(x) = B(x) - C(x) - D(x)\\}
    \onslide<3->{\text{Mercado, con impuesto } \tau\text{:}\quad & B'(x) = C'(x) + \tau\\}
    \onslide<4->{\text{Óptimo social:}\quad & B'(x^*) = C'(x^*) + D'(x^*)\\}
    \onslide<5->{\text{El mercado llega a } x^* \text{ si y solo si:}\quad & \resultado{\tau^* = D'(x^*)}}
  \end{align*}
  \onslide<6->{
  \begin{block}{Dos detalles}
    \begin{itemize}
      \item El impuesto se fija con el daño marginal \emph{en el óptimo}, no en el nivel actual de contaminación.
      \item Con $\tau=0$, el mercado iguala $B'$ con $C'$ y se pasa de $x^*$.
    \end{itemize}
  \end{block}}
  \fuente{Gruber, cap.~5.3; Salanié, cap.~10, sec.~10.1.1.}
\end{frame}

% 4 -----------------------------------------------------------------------
\begin{frame}{Específico y sobre el uso, no tasa diferenciada de IVA}
  \begin{enumerate}
    \item \textbf{El daño depende de la cantidad, no del precio.} Un litro de gasolina emite lo mismo si cuesta 20 o 25 pesos. Toca un impuesto específico, $\tau$ por litro.
    \pause
    \item \textbf{El daño es el mismo en consumo final que como insumo.} El diésel de un tráiler contamina igual que el de una camioneta familiar. Un IVA diferenciado solo alcanzaría al consumidor final: la empresa lo acredita.
  \end{enumerate}
  \pause
  \begin{alertblock}{¿Y Diamond y Mirrlees?}
    Gravar el combustible de una fábrica es la excepción a la eficiencia productiva, no su violación: corrige un precio que ya estaba mal.
  \end{alertblock}
  \pause
  \begin{block}{Orden de los impuestos}
    El IVA se aplica sobre el precio que ya incluye el impuesto específico. El específico mueve el precio relativo; el IVA, el nivel general.
  \end{block}
  \fuente{de Mooij, Hebous y Keen (2025), sec.~3.3.}
\end{frame}

% 5 -----------------------------------------------------------------------
\begin{frame}{Precio o cantidad}
  \small
  El costo de abatir puede salir alto o bajo. ¿Fijar el precio ($\tau$) o la cantidad ($\bar a$)?
  \par\smallskip
  \centering
  \begin{tikzpicture}[x=0.62cm,y=0.4cm]
    \begin{scope}
      \node at (3.2,9.2) {Daño marginal plano (CO$_2$)};
      \onslide<2->{\fill[pierde] (3.33,2.0) -- (4.17,2.5) -- (3.33,2.5) -- cycle;
      \fill[pierde] (2.5,2.5) -- (3.33,3.0) -- (3.33,2.5) -- cycle;}
      \ejes{6.4}{8.4}{abatir}{pesos}
      \draw[social] (0,2.5) -- (6,2.5);
      \node[cImpuesto,above right] at (0,2.5) {DMg $=\tau$};
      \draw[oferta] (0,0) -- (6,3.6) node[right] {bajo};
      \draw[oferta,dash pattern=on 5pt off 2pt] (0,1) -- (6,4.6) node[right] {alto};
      \node[cOferta,align=center] at (5.2,6.6) {CMg de\\abatir};
      \draw[guia] (3.33,0) -- (3.33,6.5) node[above] {$\bar a$};
    \end{scope}
    \begin{scope}[xshift=8.6cm]
      \node at (3.2,9.2) {Daño marginal empinado (umbral)};
      \onslide<2->{\fill[pierde] (3.077,1.846) -- (3.72,2.231) -- (3.72,0.56) -- cycle;
      \fill[pierde] (2.05,3.9) -- (2.05,2.231) -- (2.692,2.615) -- cycle;}
      \ejes{6.4}{8.4}{abatir}{pesos}
      \draw[social] (0,8) -- (4,0);
      \node[cImpuesto,right] at (0.3,7.9) {DMg};
      \draw[oferta] (0,0) -- (6,3.6) node[right] {bajo};
      \draw[oferta,dash pattern=on 5pt off 2pt] (0,1) -- (6,4.6) node[right] {alto};
      \draw[guia] (0,2.231) -- (6,2.231) node[right] {$\tau$};
      \draw[guia] (2.885,0) -- (2.885,6.5) node[above] {$\bar a$};
    \end{scope}
  \end{tikzpicture}
  \par
  \raggedright
  \onslide<2->{En rojo, lo que se pierde: a la izquierda con la cuota, a la derecha con el impuesto.}
  \onslide<3->{\par\textbf{Weitzman:} impuesto si el daño marginal es plano (acervo, como el CO$_2$); cuota si es empinado (umbrales).}
  \fuente{Weitzman (1974), \emph{RES} 41(4): 477--491; Gruber, cap.~5.4. Gráfica propia.}
\end{frame}

% 6 -----------------------------------------------------------------------
\begin{frame}{Con impuestos distorsionantes: recaudar y corregir}
  \begin{itemize}
    \item Hasta aquí el gobierno podía recaudar sin costo. En la realidad, cada peso cuesta más de un peso (sesión del 26 de septiembre, diap.~34).
    \pause
    \item Sandmo (1975): el impuesto óptimo a un bien que contamina tiene dos componentes que se suman.
  \end{itemize}
  \pause
  \[
    \tau_i \;=\; \underbrace{\text{componente Ramsey}}_{\text{recaudar, como en la parte 1}} \;+\; \underbrace{\text{componente pigouviano}}_{\text{corregir el daño}}
  \]
  \pause
  \begin{block}{Lectura}
    \begin{itemize}
      \item El daño solo entra en el impuesto al bien que contamina. Las demás tasas siguen la regla de Ramsey.
      \item La recaudación y la corrección son dos tareas distintas, aunque las cumpla un mismo impuesto.
    \end{itemize}
  \end{block}
  \fuente{Sandmo (1975), \emph{Swedish Journal of Economics} 77(1): 86--98; de Mooij, Hebous y Keen (2025), nota 84; Salanié, sec.~10.1.2.}
\end{frame}

% 7 -----------------------------------------------------------------------
\begin{frame}{¿Doble dividendo?}
  \begin{columns}[T]
    \begin{column}{0.48\textwidth}
      \begin{block}{Versión débil}
        Usar la recaudación verde para bajar otro impuesto distorsionante es mejor que devolverla de suma fija.\\
        \emph{Casi nadie la discute.}
      \end{block}
      \begin{block}{Versión fuerte}
        Además de limpiar el aire, el cambio mejora la eficiencia del resto del sistema fiscal: «sale gratis».
      \end{block}
    \end{column}
    \begin{column}{0.48\textwidth}
      \pause
      \begin{alertblock}{Bovenberg y de Mooij (1994)}
        \begin{itemize}
          \item El impuesto verde sube los precios y baja el salario real: agrava la distorsión en el mercado de trabajo (\textbf{efecto interacción}).
          \item Ese efecto suele ser mayor que la ganancia por reciclar la recaudación.
          \item El impuesto ambiental óptimo queda \textbf{por debajo} del pigouviano.
        \end{itemize}
      \end{alertblock}
      \pause
      {\small De Mooij es coautor del artículo del FMI de la parte 1.}
    \end{column}
  \end{columns}
  \fuente{Bovenberg y de Mooij (1994), \emph{AER} 84(4): 1085--1089; Salanié, sec.~10.2.}
\end{frame}

% 8 -----------------------------------------------------------------------
\begin{frame}{¿Quién paga el impuesto verde?}
  \begin{itemize}
    \item En los países ricos, un impuesto al carbono suele ser \textbf{regresivo}: todos usan energía, y los pobres gastan en ella una parte mayor de su ingreso.
    \pause
    \item En países de ingreso bajo y medio puede ser \textbf{progresivo}: el auto y el combustible se concentran arriba. Dorband et al. (2019) lo encuentran, en promedio, para países con ingreso per cápita menor a 15 mil dólares (PPA de 2011).
    \pause
    \item Un metaanálisis de 53 estudios: es más probable que resulte progresivo en países de menor ingreso y en transporte.
  \end{itemize}
  \pause
  \begin{block}{Lo que decide}
    El uso de la recaudación. Si se devuelve como una transferencia igual por persona, el hogar pobre recibe, en proporción a su ingreso, mucho más de lo que paga el rico.
  \end{block}
  \fuente{Dorband, Jakob, Kalkuhl y Steckel (2019), \emph{World Development} 115: 246--257; Ohlendorf, Jakob, Minx, Schröder y Steckel (2021), \emph{Environmental and Resource Economics} 78(1): 1--42.}
\end{frame}

% 9 -----------------------------------------------------------------------
\begin{frame}{Internalidades: el daño que uno se hace a sí mismo}
  \begin{itemize}
    \item Una \textbf{internalidad} es un costo que el consumidor no toma en cuenta al decidir, aunque lo pague él: el fumador que mañana sí va a dejarlo.
    \item El IEPS a bebidas azucaradas: en 2026, 3.0818 pesos por litro con azúcares añadidos y 1.50 con edulcorantes (sesión del 26 de septiembre, diap.~35--37).
  \end{itemize}
  \pause
  \begin{block}{Aquí el caso del impuesto específico es menos claro}
    \begin{itemize}
      \item El daño del CO$_2$ es por litro de combustible. El del azúcar depende de quién bebe y cuánto: no es igual para todos los litros.
      \item Y el que paga el impuesto es el mismo que se beneficia de la corrección.
    \end{itemize}
  \end{block}
  \fuente{de Mooij, Hebous y Keen (2025), sec.~3.3; Gruber, cap.~6; LIEPS, art.~2, fr.~I, inciso G (reforma DOF 07-11-2025).}
\end{frame}

% =========================================================================
% II. MÉXICO
% =========================================================================
\section{México}

% 10 ----------------------------------------------------------------------
\begin{frame}{IEPS a gasolinas y diésel: ¿pigouviano o recaudatorio?}
  \begin{columns}[T]
    \begin{column}{0.46\textwidth}
      \centering
      \begin{tabular}{lr}
        \toprule
        Cuota 2026 & \$ por litro\\
        \midrule
        Gasolina, menos de 91 oct. & 6.7001\\
        Gasolina, 91 oct. o más & 5.6579\\
        Diésel & 7.3634\\
        \bottomrule
      \end{tabular}
      \par\smallskip
      \raggedright
      {\small Cuota fija por litro; se actualiza cada año por inflación.}
    \end{column}
    \begin{column}{0.5\textwidth}
      \pause
      \begin{itemize}
        \item Cada semana Hacienda publica un \textbf{estímulo}: un porcentaje de la cuota que no se cobra, para suavizar el precio.
        \item En varias semanas de 2022 el estímulo llegó a 100\,\%: la cuota quedó en cero y hubo estímulos adicionales.
      \end{itemize}
    \end{column}
  \end{columns}
  \pause
  \begin{alertblock}{Un estímulo es un subsidio que va contra Pigou}
    Abarata justo lo que contamina. Y como el consumo de gasolina se concentra en los hogares de mayor ingreso, el subsidio también.
  \end{alertblock}
  \fuente{LIEPS, art.~2, fr.~I, inciso D; Acuerdo 179/2025 de actualización de cuotas (DOF, 22-12-2025); acuerdos semanales de estímulos fiscales (DOF, 2022).}
\end{frame}

% 11 ----------------------------------------------------------------------
\begin{frame}{IEPS a combustibles fósiles: el «impuesto al carbono» de 2014}
  \begin{columns}[T]
    \begin{column}{0.46\textwidth}
      \centering
      \begin{tabular}{lr}
        \toprule
        Cuota 2026 & \\
        \midrule
        Gasolinas & 17.76 ¢/litro\\
        Diésel & 21.55 ¢/litro\\
        Combustóleo & 23.00 ¢/litro\\
        Carbón mineral & 47.12 \$/tonelada\\
        \bottomrule
      \end{tabular}
      \par\smallskip
      \raggedright
      {\small La cuota sigue el contenido de carbono de cada combustible.}
    \end{column}
    \begin{column}{0.5\textwidth}
      \pause
      \begin{block}{¿Cuánto vale una tonelada de CO$_2$?}
        \begin{itemize}
          \item Precio implícito del impuesto mexicano: 3.9 dólares por tonelada (Banco Mundial, abril de 2025). Cubre 29\,\% de las emisiones.
          \item Costo social del carbono de la EPA de Estados Unidos: 190 dólares por tonelada.
        \end{itemize}
      \end{block}
      \pause
      Cobramos unas 50 veces menos que el daño estimado.
    \end{column}
  \end{columns}
  \fuente{LIEPS, art.~2, fr.~I, inciso H; Acuerdo 179/2025 (DOF, 22-12-2025); Banco Mundial (2025), \emph{State and Trends of Carbon Pricing}; US EPA (2023), \emph{Report on the Social Cost of Greenhouse Gases} (dólares de 2020, emisiones de 2020, descuento de 2\,\%).}
\end{frame}

% 12 ----------------------------------------------------------------------
\begin{frame}{Impuestos verdes estatales: muchos en la ley, pocos en la caja}
  \begin{columns}[T]
    \begin{column}{0.44\textwidth}
      \begin{itemize}
        \item 26 entidades los contemplan en su ley de ingresos; solo 13 esperan recaudar algo.
        \item 2,590.7 millones de pesos en 2023: en promedio, 3.5\,\% de sus ingresos por impuestos.
      \end{itemize}
    \end{column}
    \begin{column}{0.54\textwidth}
      \centering
      \begin{tikzpicture}
        \begin{axis}[
          width=5.6cm, height=4.6cm,
          xbar, bar width=10pt,
          xmin=0, xmax=55,
          xlabel={\% del total},
          xlabel style={font=\small},
          symbolic y coords={Residuos,Suelo y agua,Emisiones al aire,Extracción},
          ytick=data,
          y tick label style={font=\small},
          x tick label style={font=\small},
          nodes near coords,
          every node near coord/.append style={font=\small},
          axis lines*=left,
          enlarge y limits=0.18,
        ]
          \addplot[fill=cRecauda!50,draw=cRecauda] coordinates {(0.6,Residuos) (25.9,Suelo y agua) (30.4,Emisiones al aire) (43.2,Extracción)};
        \end{axis}
      \end{tikzpicture}
    \end{column}
  \end{columns}
  \pause
  \begin{block}{Lectura}
    Lo que más recauda es la extracción de materiales, no las emisiones. ¿Corrige un daño o es una base fácil de cobrar?
  \end{block}
  \fuente{CIEP, Sánchez Salazar (2023), «Impuestos verdes: una oportunidad para fortalecer la recaudación subnacional», 25-10-2023. Datos de las leyes de ingresos estatales de 2023.}
\end{frame}

% 13 ----------------------------------------------------------------------
\begin{frame}{Para discutir: ¿qué está bien y qué no?}
  Dos afirmaciones de la nota del CIEP:
  \begin{enumerate}
    \item «La hipótesis del doble dividendo, planteada por Pearce (1991), sirve para justificar un aumento en la recaudación por medio de una reforma fiscal verde.»
    \item «Es necesario que se grave a la producción final, puesto que gravar a productos intermedios aumenta los costos.»
  \end{enumerate}
  \pause
  \begin{itemize}
    \item Contrasten con la diapositiva 7 (Bovenberg y de Mooij) y con la diapositiva 4 (el daño es el mismo como insumo).
    \item Una idea de la nota que sí se sostiene: rediseñar la tenencia según el tamaño y las características del motor.
  \end{itemize}
  \pause
  \begin{alertblock}{Pregunta}
    ¿Cuál de las dos afirmaciones defenderían y cuál corregirían? ¿Por qué?
  \end{alertblock}
  \fuente{CIEP (2023), citado arriba; Pearce (1991), \emph{Economic Journal} 101(407): 938--948.}
\end{frame}

% =========================================================================
% III. CIERRE
% =========================================================================
\section{Cierre}

% 14 ----------------------------------------------------------------------
\begin{frame}{Ideas centrales: impuestos verdes}
  \begin{enumerate}
    \item Corregir pide un impuesto específico, sobre la emisión, donde ocurra: en el consumo o en la producción.
    \item Precio o cantidad depende de la forma del daño marginal.
    \item Con impuestos distorsionantes, el óptimo queda por debajo de Pigou. El doble dividendo fuerte rara vez aparece.
    \item Quién paga depende de qué se hace con lo recaudado.
  \end{enumerate}
\end{frame}

% 15 ----------------------------------------------------------------------
\begin{frame}{Cierre de la sesión}
  \begin{columns}[T]
    \begin{column}{0.32\textwidth}
      \begin{block}{El IVA recauda}
        Base amplia, tasa uniforme, sin gravar insumos.
      \end{block}
    \end{column}
    \begin{column}{0.32\textwidth}
      \begin{block}{Los verdes corrigen}
        Impuestos específicos, aun sobre insumos.
      \end{block}
    \end{column}
    \begin{column}{0.32\textwidth}
      \begin{block}{Las transferencias redistribuyen}
        Mejor que las tasas reducidas.
      \end{block}
    \end{column}
  \end{columns}
  \vspace{1em}
  \begin{alertblock}{Primer parcial: 20 de octubre}
    De teoría. La parte de tasa uniforme (Ramsey, Atkinson y Stiglitz) entra.
  \end{alertblock}
  \begin{block}{Próxima sesión: 27 de octubre}
    Un marco unitario para el sistema fiscal (Kaplow) y la digitalización. El pago fraccionado y la factura electrónica son el puente.
  \end{block}
\end{frame}

\end{document}
